\scanpage{86}
% L1
\[
\Theta(v\mid t)=-\frac{1}{\unsure{i|1}}\sum_{-\infty}^{+\infty}(-1)^n e^{(n-\frac12)^2i\pi t+(2n-1)i\pi v}
\tag*{\textcircled{\scriptsize 1}}
\]

% L2–4
série qui peut aussi s'écrire, en \deleted{\unsure{groupant}} décomposant la somme en $\sum_1^\infty$ et $\sum_{-\infty}^{\unsure{0|o}}$ et groupant les termes deux à deux

% L5
\[
\Theta(v\mid t)=2_{\deleted{\unsure{i|1}}}\sum_1^\infty(-1)^{n-1}e^{(n-\frac{1}{\unsure{1|2}})^2i\pi t}\sin(2n-1)\pi v
\tag*{\textcircled{\scriptsize 2}}
\]

% L6–7
On a alors, par construction. $\Theta$ est définie \inserted{et holomorphe} pour $\unsure{\mathfrak I|J}t>0$ et pour $v$ quelconque

% L8
\[
\sigma(u\mid\omega_1,\omega_2)=
\frac{2\omega_1}{\Theta'_v(0\mid\frac{\omega_2}{\omega_1})}\,
e^{\frac{\eta_1u^2}{2\omega_1}}\,\Theta\left(\frac{u}{2\omega_1},\frac{\omega_2}{\omega_1}\right)
\tag*{\textcircled{\scriptsize 3}}
\]

% L9–11
La fonction $\Theta$ est exprimée par une série très rapidement convergente. -- $\Theta(v\mid t)$ \unsure{est} fonction impaire de $v$, et on a

% L12–13
\[
\left\{\begin{aligned}
\Theta(v+1\mid t)&=-\Theta(v\mid t)&\Theta(v+2\mid t)&=\Theta(v\mid t)\\
\Theta(v+t\mid t)&=\Theta(v\mid t)e^{-2i\pi v-i\pi t-i\pi}
\end{aligned}\right.
\tag*{\textcircled{\scriptsize 4}}
\]

% L14–16
Ces relations caractérisent $\Theta$ \unsure{a|à} un facteur constant près, pour $t$ donné (comme on voit par le calcul qui avait donné $\Theta$)

% L17
Notons aussi la formule (vérification immédiate)

% L18
\[
\Theta(v\mid t+4)=-\Theta(v\mid t),\qquad\text{d'où}\quad\Theta(v\mid t+8)=\Theta(v\mid t)
\tag*{\textcircled{\scriptsize 5}}
\]

% L19–20
Les zéros de $\Theta(v\mid t)$ sont déterminés par ceux de $\sigma u$, ce sont les nombres $v=m_1+m_2 t$. --

% L21–22, foot
Pour une expression de \textcircled{\scriptsize 1}, \textcircled{\scriptsize 2}, commode au calcul, on pose parfois $h(t)=h=e^{i\pi t}$; alors

% lower left
% archive sheet number: 85
